%% kapa: exam=1 round=20 year=2009 period=2 form=A subject=appraisal_law subject_ko=부동산관계법규 seq=27 number=67
%% source: Q-Net 감정평가사 기출문제 (cst00309) artlSeq=5024178 file=extracted/1008/5024178/2009년 제20회 감정평가사 1차시험 기출문제/2교시.hwp
%% conversion: kordoc 4.24.51 -> tools/gen1.py (auto); review=proofread-sonnet; flags=none
\begin{question}{67}
국유재산법령상 행정재산의 사용ㆍ수익허가에 관한 설명으로 옳지 \uline{않은} 것은?
\begin{choices}
  \item 행정재산은 그 용도 또는 목적에 장애가 되지 아니하는 범위 안에서 그 사용ㆍ수익을 허가할 수 있다.
  \item 사용ㆍ수익의 허가를 받은 자는 관리청이 그 행정 또는 보존 목적의 수행에 필요하다고 인정하는 시설물의 기부를 전제로 한 축조를 제외하고는 그 허가받은 재산상에 건물 기타의 영구시설물을 축조하지 못한다.
  \item 기부를 채납한 재산에 대하여 사용ㆍ수익 허가를 받은 자가 당해 재산의 기부자인 경우에는 관리청의 승인을 얻어 다른 사람으로 하여금 당해 재산을 사용ㆍ수익하게 할 수 있다.
  \item 관리청이 행정재산의 사용허가를 취소하거나 철회하려는 경우에 청문의 실시여부는 관리청의 재량사항이다.
  \item 관리청은 행정재산의 사용ㆍ수익 허가를 받은 자가 당해 재산의 보관을 해태하였거나 그 사용목적에 위배한 때에는 그 허가를 취소 또는 철회할 수 있다.
\end{choices}
\end{question}
\answer{4}
